\documentclass[10pt,a4paper]{amsart}
\usepackage[margin=19mm]{geometry}
\usepackage{amsmath,amssymb,mathtools}
\usepackage[scr=rsfs,cal=euler]{mathalfa}
\usepackage{xfrac}
\usepackage[unicode,hidelinks,bookmarks=false]{hyperref}
\newcommand{\ie}{i.e.\ }
\newcommand{\cf}{cf.\ }
\newcommand{\resp}{respectively\ }
\newcommand{\Subsection}{\subsection}
\providecommand{\nobreakdash}{\nobreak\hbox{-}}
\setlength{\emergencystretch}{2em}
\multlinegap0pt
\usepackage{amssymb}
\usepackage[all]{xy}
\usepackage{enumerate}
\usepackage{dsfont}
\usepackage{bm}
\newtheorem{lemma}{Lemma}
\newtheorem{definition}{Definition}
\newcommand{\Avnat}{\Linfty(\mathbb{R}) \overline{\otimes} \mathcal{M}}
\newcommand{\C}{\mathbb{C}}
\newcommand{\Linfty}{L_\infty}
\newcommand{\Lone}{L_1}
\newcommand{\Ltwo}{L_2}
\newcommand{\R}{\mathbb{R}}
\newcommand{\id}{\mathrm{Id}}
\newcommand{\la}{\langle}
\title{Calder\'on-Zygmund theory with noncommuting kernels via $H_1^c$}
\author{Antonio Ismael Cano-M\'armol, \'Eric Ricard}
\date{Source: arXiv:2309.01266v3 (CC BY 4.0)}
\begin{document}
\maketitle

\noindent\emph{Source and licence.} Calder\'on-Zygmund theory with noncommuting kernels via $H_1^c$. Version arXiv:2309.01266v3 (CC BY 4.0), distributed by arXiv under the Creative Commons Attribution 4.0 International licence, http://creativecommons.org/licenses/by/4.0/, as recorded in the arXiv licence metadata for this exact version and shown on the abstract page https://arxiv.org/abs/2309.01266v3. Full licence notice: \texttt{licenses/arxiv-2309.01266-CC-BY-4.0.txt}.\par\smallskip

\noindent\textit{Editorial scope.} Counted unit: the article's lemma lem:approxKernel (source0.tex lines 1272--1287). The article's complete original proof of it is delivered with it (source0.tex lines 1290--1354, one \texttt{proof} environment of 65 lines). No mathematics is written, repaired or re-derived: the statement and the proof are the article's own lines, in the article's own notation, and no retained line is substituted or re-flowed.  Retained uncounted as the source-local context the proof needs: lines 1020--1029; lines 1178--1207.  Nothing was omitted from any retained span, so the package carries no omission record.  No retained line contains a citation, so the wrapper carries no bibliography and no bibliography entry.  No retained line contains a figure environment, an image-inclusion command or drawing code, so no image is delivered and the package declares no asset dependency.  Editorial formatting only, in addition: an AMS \texttt{amsart} preamble that restates the article's own declarations and macros used by the retained lines, namely the environments \texttt{lemma}, \texttt{definition} on separate counters, together with the standard packages those lines need.  The retained environments print in this document as Lemma 1, Definition 1 in order of appearance, whereas the article prints the same items with the numbering of its own full document.  The complete native source of the article as distributed in the arXiv e-print for this exact version is bundled unchanged as \texttt{sources/146483/source0.tex}.
\begin{lemma}\label{lem:convolutionAtoms}
        Given a $c$-atom $a=bh$ supported on $I$, the following holds:
        \begin{enumerate}
        \item   $\frac 12\phi_n*a$ is a finite convex combination of $c$-atoms
        \item If $n\geq \frac 2{|I|}$, $\phi_n* a - a= \lambda_n a_n$,
          where $a_n$ is a $c$-atom supported on $2I$,  $\lambda_n\in \C$ with
          $\lim_{n\to \infty} \lambda_n=0$.
        \item We have $\gamma(\phi_n*a)=C_{\phi_n}(a)$ for all $n\geq 1$.
                  \end{enumerate}
      \end{lemma}

    \begin{definition}\label{def:CZ}
        Assume that $T$ is a bounded operator on $\Ltwo(\Linfty(\mathbb{R}) {\overline{\otimes}} \mathcal{M})$. We say that $T$ is a \emph{Calder\'on-Zygmund operator} if there exists some function
        \begin{align*}
            K : \mathbb{R} \times \mathbb{R} \ \setminus \ \{x=y\} \longrightarrow \mathcal{M}
        \end{align*}
        such that for every pair of intervals $I,J$ satisfying $d(I,J) > 0$, there exist $K_{I,J} \in \Linfty(I \times J) \overline{\otimes} \mathcal{M}$ such that
        \begin{itemize}
          \item ${K}_{I,J}(t) = K(t)$ for  almost every $t \in I \times J$,
 \item   for any $f,g \in \mathcal S$ supported respectively on $J$ and $I$:  \begin{align}\label{dissup}
            \Big(\tau \circ \int\Big)(g \ Tf) = \langle f(y)g(x), K_{I,J} \rangle_{\Lone(\Linfty(I\times J)\overline{\otimes} \mathcal{M}),\Linfty(I \times J)\overline{\otimes} \mathcal{M}}
        \end{align}
      \end{itemize}
      
        We say that $T$ is a \emph{left Calder\'on-Zygmund operator} if it also  satisfies that for any compactly supported $f \in \Ltwo(\Avnat)$ and $h \in \mathcal{M}$:
        \begin{align*}
            T(fh) = T(f)h. 
        \end{align*}

      
       Let $\lambda>0$, we say that the kernel $K$ (or $T$) satisfies the \emph{Hörmander condition} for $\lambda$ whenever
        \begin{align}\label{eq:hormanderCondition}
            \int_{|x-y| \geq \lambda |y^\prime-y|} \|K(x,y) - K(x,y^\prime)\|_{\mathcal{M}} \ dx \leq C_\lambda
        \end{align}
        for some constant $C_\lambda \geq 0$.  Regard that should the Hörmander condition hold for $\lambda > 0$, then it holds for $\lambda' > \lambda$.
        % On the other hand, $K$ will be said to satify the \emph{size condition} whenever
        % \begin{align}\label{eq:sizeCondition}
        %     \|K(x,y)\|_{\mathcal{M}} \leq \frac{C^\prime}{|x-y|} 
        % \end{align}
        % for $x \neq y$.
    \end{definition}

    \begin{lemma}\label{lem:approxKernel}
      Let $T$ be a left Calder\'on-Zygmund operator whose associated
      kernel $K$ satisfies the Hörmander condition  for some $\lambda \geq 1$\eqref{eq:hormanderCondition}. Then, there exist a sequence of left Calder\'on-Zygmund operators $(T_m)_{m\geq 1}$ with
      kernels
      $(K_m)_{m \geq 1} \subseteq \Linfty(\mathbb{R} \times
      \mathbb{R}) \overline{\otimes} \mathcal{M}$ and a constant
      $C'>0$ such that
        \begin{align}\label{eq:weakerHormander}
            \int_{|x-y| \geq 2\lambda |y^\prime-y|} \|K_m(x,y) - K_m(x,y^\prime)\|_{\mathcal{M}} \ dx \leq C'.
        \end{align}
         Moreover, there holds
            \begin{align*}
                \lim_{m \rightarrow \infty} \|T_m(f) - T(f) \|_{\Ltwo(\mathbb{R};\Ltwo(\mathcal{M}))} = 0
            \end{align*}
            for every $f \in \Ltwo(\Avnat)$.
    \end{lemma}

    \begin{proof} 
    We consider an even smooth function $\phi$ which is supported on $(-1,1)$ with
      $\int \phi = 1$ as in Lemma ~\ref{lem:convolutionAtoms}. Then, set $R_m$ to be the operator defined on
      $\Ltwo(\mathbb{R};\Ltwo(\mathcal{M}))$ as $R_m(f) = f * \phi_m$
      where $\phi_m(x) = m \phi(mx)$.
    
      We define the function $K_m:\R^2\to \mathcal M$ by duality. Assume
        $z\in \Lone(\mathcal M)$ is decomposed as $z=ab$ with $a,b\in \Ltwo(\mathcal M)$, and define the quantity 
      \begin{align}\label{formulaKm}
        \la K_m(x,y),z\rangle_{\mathcal M,\Lone(\mathcal M)}  =
        \int \tau\big(T((\tau_y\phi_m)(t)a)\tau_x ({\phi}_m)(t)b\big) dt,
      \end{align}
      where $\tau_x \phi(t) = \phi(t-x)$. Indeed assume $z=\alpha\beta$ with $\alpha,\beta\in \Ltwo(\mathcal M)$ and choose $p$ a
      finite projection so that $pa, p\alpha\in \mathcal M$, then using the right-modularity condition,
      $$\int_\R \tau\big(T(\tau_y\phi_mpa)\tau_x{\phi}_m \ b\big)= \int_\R \tau\big(T(\tau_y\phi_m p)\tau_x \phi_m \ z\big)=\int_\R \tau\big(T(\tau_y\phi_m p\alpha)\tau_x\phi_m \ \beta\big).$$
      Letting $p\to 1$ and using the $L_2$-continuity of $T$ yields
      $$\langle T(\tau_y\phi_ma), \tau_x ({\phi}_m)b \rangle_{\Ltwo(\R;\Ltwo(\mathcal M))}=\langle T(\tau_y\phi_m\alpha), \tau_x ({\phi}_m)\beta \rangle_{\Ltwo(\R;\Ltwo(\mathcal M))}.$$      
      Linearity is proved in a similar way. We also get
      \begin{align*}
            \|K_m(x,y)\|_{\mathcal{M}} &= \sup_{\substack{{a,b\in \Ltwo(\mathcal M)}\\ \|a\|_2=\|b\|_2=1}} | (\int \circ \tau) \big(T(\tau_y \phi_ma)  \ \tau_x {\phi}_m  b \big) | \\
            &\leq \|T\|_{\Ltwo(\R;\Ltwo(\mathcal M))}  \| \phi_m\|_{\Ltwo(\mathbb{R})}^2.
        \end{align*}
        Assuming that $f$ and $g$ are simple functions, we get by linearity 
\begin{align*}
  \langle R_m T R_m f , g \rangle &= \tau \int \int T\Big(\int \phi_m(\cdot-y) \ f(y) \ dy\Big)(t) \ \phi_m(t-x) \ dt \ g(x) \ dx\\
  &= \int\int \tau\big(K_m(x,y)f(y)g(x)\big) \ dx \ dy.
\end{align*}
The formula extends to $f,g\in \Ltwo(\R;\Ltwo(\mathcal M))$ by continuity since $K_m$ is bounded on the whole $\R^2$. Therefore, $K_m$ is the kernel for $R_m T R_m$ in the sense of Definition~\ref{def:CZ}. Moreover, $K_m$ is a continuous function $\R\times\R \to \Linfty(\mathcal M)$ by continuity of translations on $\Ltwo(\R)$ as
\begin{align}
  \nonumber \| K_m(x,y)-K_m(x,y')\|_\mathcal M&\leq                                           \|T\|  \ \| \phi_m\|_{2} \ \| (\tau_y-\tau_{y'})\phi_m\|_2 \\ \nonumber 
  &\leq \|T\| \ m^{1/2} \ \|\phi\|_2 \ \big( 2 m^{3/2}  {|y-y'|} \|\phi'\|_{\infty}\big)\\ \label{lipsest} & =
2  m^2 {|y-y'|} \|T\| \|\phi\|_2 \|\phi'\|_{\infty}
\end{align}
A similar Lipschitz estimate holds for $\| K_m(x,y)-K_m(x',y)\|_\mathcal M$. 
      
On the other hand, the kernel $K_m$ satisfies the Hörmander condition for $2\lambda$. First, consider $x,y,y'$ such that $|x-y| \geq 2\lambda |y^\prime - y|$
and $|x-y| > \frac{4}{m}$, then $|x-y'|\geq \frac{|x-y|}2>\frac 2 m$ as $\lambda\geq 1$. Thus the support of $\tau_x {\phi}_m$ is
disjoint from that of $\tau_y \phi_m$ and $\tau_{y'} \phi_m$. Let
$I_{y,y^\prime}$ and $I_x$ denote some disjoint closed intervals containing
the support of $\tau_y \phi_m,\tau_{y^\prime}\phi_m$ and
$\tau_x {\phi}_m$ respectively. Then, the identities \eqref{dissup} and \eqref{formulaKm} for the kernel
$K_{I_x,I_{y,y^\prime}}$ (extending to 0 outside $I_x\times I_{y,y^\prime}$) give
$$ K_m(x,z)= K_{I_x,I_{y,y^\prime}}*({\phi}_m\otimes \phi_m)(x,z)$$
for $z=y,y'$. Thus we can write
        \begin{align*}
            &\int_{|x-y| \geq 2\lambda |y^\prime - y|\vee \frac{4}{ m} } \|K_m(x,y)-K_m(x,y^\prime)\|_{\mathcal{M}} \ dx \\
                        % &= \int_{|x-y| \geq 4 |y^\prime -y|} \sup_g \big| \tau \int \int \big[K_{I_x,I_{y,y^\prime}}(x - \frac{u}{m},y + \frac{v}{m}) - K_{I_x,I_{y,y^\prime}}(x - \frac{u}{m},y^\prime + \frac{v}{m}) \big]  \ \phi(u) \phi(v) g \ du \ dv \big| \ dx \\
            &\leq\int_{|x-y| \geq 2\lambda |y^\prime -y|\vee \frac{4}{m} } \int \int \big\| {K}(x - u,y -{v}) - {K}(x - {u},y^\prime - {v}) \big\|_{\mathcal{M}} \ {\phi_m}(u) \phi_m(v)\ du \ dv \ dx.
        \end{align*}
        Using that for $|u|,|v|\leq \frac{1}{m}$ (i.e. in the support of $\phi_m$), 
$$|x - u - (y -{v}) | \geq |x-y| - \frac{2}{m} > \frac{1}{2} |x-y| \geq \lambda |y^\prime-v -(y-v)|,$$
the H\"ormander condition for the kernel $K$ implies
    $$\int_{|x-y| \geq 2\lambda |y^\prime - y|\vee \frac{4}{ m}} \|K_m(x,y)-K_m(x,y^\prime)\|_{\mathcal{M}} \ dx \leq C.$$ 
    On the other hand, whenever $|x -y | \leq 4/ m$ holds, the estimate \eqref{lipsest} yields
        \begin{align*}
            &\int_{\frac{4}{m} \geq |x-y| \geq 2\lambda |y^\prime -y|} \|K_m(x,y)-K_m(x,y^\prime)\|_{\mathcal{M}} \ dx \\
            &\leq 2 m^2 \|T\| \|\phi^\prime\|_{\infty}\|\phi\|_2 \int_{\frac{4}{m} \geq |x-y| \geq 2\lambda |y^\prime -y|} |y^\prime -y| \ dx \leq \frac{8}{\lambda} \ \|T\| \|\phi^\prime\|_{\infty} \|\phi\|_2.
        \end{align*}
        Thus $K_m$ satisfies condition \eqref{eq:weakerHormander}. Finally, given $f \in \Ltwo(\Avnat)$ there holds
        \begin{align*}
            \|T_m(f) - T(f)\|_{\Ltwo} &\leq \|R_m T (R_m(f) - f)\|_{\Ltwo} + \|(R_m - \id)T(f)\|_{\Ltwo} \\
            &\leq \|T\| \ \|R_m(f) - f\|_{\Ltwo} + \|(R_m - \id)T(f)\|_{\Ltwo} \rightarrow 0.
        \end{align*}
        as $m$ grows to infinity.
  \end{proof}

\end{document}
